\chapter{A Machine Made of Living Neurons}
\chaptermark{A Machine Made of Living Neurons}
\label{sec:livingmachine}
\begin{chaptergoals}
\item why a culture with no broadcast channels fires network bursts, and what sets their interval;
\item how living networks change their responses under feedback, and what remains unproven;
\item how an organoid serves as a reservoir whose linear readout is trained outside, in silicon;
\item how cultures have been taught so far, and why the teacher still stands outside the dish.
\end{chaptergoals}

\section{A test bench with an open schematic}

In Chapter~\ref{sec:debugging} we debugged a machine that has no schematic: a probe hears a thousand neurons out of eighty billion, and everything else has to be guessed. An engineer in that position would do something simple: build a small piece of the circuit on a breadboard, where every part is visible. With neurons, this can be done too.

Cortical neurons are dissociated and grown on a chip with an electrode array, from a few dozen to tens of thousands of electrodes. Within a few weeks the cells sprout processes and wire up into a network of thousands or hundreds of thousands of neurons, mostly \nt{GLUT} with a smaller fraction of \nt{GABA} that depends on the preparation: in hippocampal cultures it is about $6\%$ of the neurons~\cite{Benson1994}. Every electrode can both listen, like the probe of Chapter~\ref{sec:debugging}, and inject current, like a test signal. The second kind of bench uses \emph{organoids}: three-dimensional clumps of neural tissue about a millimeter across, grown from human stem cells~\cite{Lancaster2013}. On such benches living networks run for months; there are facilities where experiments on organoids can be run remotely, over the internet, for more than a hundred days in a row~\cite{Jordan2024}. The field is called \emph{biological computing} or \emph{organoid intelligence}~\cite{Smirnova2023}.

The bench differs from the brain in two important ways. It is tiny: it has a hundred thousand times fewer neurons. And it has no broadcast channels: a cortical culture contains no \nt{DOPA}, \nt{SERO}, \nt{NORA}, or \nt{ACET} neurons. It is a machine with a data bus and flow control but no control registers. Let us see what it can do.

\section{What the network does on its own}

Left alone, a culture neither stays silent nor produces uniform noise. From time to time almost the whole network fires at once, in a \emph{network burst}, and then falls quiet again; the intervals between network bursts range from a second to minutes, depending on the culture~\cite{Wagenaar2006}. To an engineer this is a familiar picture: an amplifier with strong positive feedback and no load turns into an oscillator.

We already know the mechanism. Strong mutual connections between \nt{GLUT} neurons set off an avalanche, as in Chapter~\ref{sec:gaba} when the conditions of Proposition~\ref{prop:ei} are violated. The avalanche quickly depletes the transmitter store in the synapses~\eqref{eq:depression}, and the network goes silent. The store recovers with time constant $\tau_{\text{rec}}$, and once the fraction $x^*$ needed for a new avalanche has recovered, a new network burst follows. The interval between network bursts is
\begin{empheq}[box=\keybox]{equation}
T_{\text{burst}} \;\approx\; \tau_{\text{rec}}\,\ln\frac{1}{1-x^*} .
\label{eq:burstinterval}
\end{empheq}
With $\tau_{\text{rec}}=0.8\,\text{s}$ from Chapter~\ref{sec:code} and $x^*=0.9$ this gives about two seconds; with $x^*=0.99$, about four. This is a model estimate with the book's $\tau_{\text{rec}}$, and it lands at the short end of the measured intervals. Direct measurement in microcultures shows the transmitter store recovering over tens of seconds after a network burst~\cite{CohenSegal2011}, so $\tau_{\text{rec}}$ is larger there; with such a $\tau_{\text{rec}}$, formula~\eqref{eq:burstinterval} gives tens of seconds to minutes, as in slow cultures. This is a relaxation oscillator, a relative of the rhythm generator of Chapter~\ref{sec:muscles}: there it was recharged by the fatigue of neuron groups, here by the transmitter store.

Network bursts are what the network does when it has nothing to do. In the living brain a similar picture appears in deep sleep (Chapter~\ref{sec:sleep}) and in the developing brain, before real inputs arrive. The resemblance is suggestive, but that the mechanisms are the same is a hypothesis.

\section{Teaching a network without a dopamine teacher}

Since the culture has no \nt{DOPA}, the ``better or worse'' signal has to be supplied by us, through the electrodes. Experiments show that a network on the bench really does change its responses, and the most interesting thing about them is \emph{what} it is taught with.

\emph{A teacher who falls silent.} The network is stimulated through one electrode once every one to three seconds until the desired response appears on another electrode $50\pm10\ms$ after the stimulus. As soon as it appears, stimulation is switched off for five minutes. After several such cycles the desired response appears right away~\cite{Shahaf2001}. The reward here is silence: the stimulus shakes the network, and ending the stimulus leaves it in the state that produced the right answer. This is gradient descent by random perturbation from Chapter~\ref{sec:dofa} (Proposition~\ref{prop:gradient}), where the shaking itself plays the role of the error: keep shaking until it comes out right.

\emph{A network that plays tennis.} In an experiment on a bench with a dense electrode array, about a million neurons were connected to a computer game of one-paddle tennis~\cite{Kagan2022}. The ball position was delivered as current to ``sensory'' electrodes: where the ball is, by a place code; how far away, by frequency. Activity in two ``motor'' regions moved the paddle up or down. The learning rule was unusual. If the paddle hit the ball, the network got a short \emph{predictable} stimulus; if it missed, it got several seconds of \emph{unpredictable} random current. Over twenty minutes of play, runs of hits grew longer. A significant improvement within a session appeared only with full feedback; with silence in place of the stimulus the culture also played better by the end of the session than without feedback, but with a smaller effect, and there was no robust learning from session to session over several days. The experiment is analyzed in detail in Chapter~\ref{sec:dishbrain}.

The authors explained the result with the hypothesis that the network acts so as to make its input predictable~\cite{Friston2010}. This is the same idea as spike economy in Chapter~\ref{sec:code} and frame prediction in Chapter~\ref{sec:recognition}: the machine spends less when the input is expected. But both the experiment and, especially, the authors' talk of the culture's ``sentience'' drew sharp criticism: the effect is small, and it can be explained more simply~\cite{Balci2023}. An honest assessment is this: the culture on the bench changes its behavior under feedback, which is measured; that it learns specifically to predict its input is a hypothesis.

\emph{A network that drives a body.} In a similar way, a culture was connected to a simple virtual ``body'' and taught to move it toward a target by choosing the spatiotemporal pattern of training stimuli~\cite{Bakkum2008}. None of these experiments involves dopamine: the teacher is a current pattern chosen by the engineer. What changes when the teacher becomes a program, or dopamine itself, is the subject of Sections~\ref{sec:dishdopamine}--\ref{sec:cartpole}.

\begin{figure}[t]
\centering
\begin{tikzpicture}[>=Stealth,
  blk/.style={draw,very thick,rounded corners=2pt,align=center,font=\footnotesize,minimum height=0.8cm,fill=blue!5}]
% (a) closed loop
\node[font=\small] at (1.5,4.3) {(a) closed loop};
\node[blk,text width=2.6cm] (G) at (1.5,3.3) {game: ball and paddle};
\draw[very thick,fill=orange!10,rounded corners=3pt] (0.5,0.2) rectangle (2.5,1.6);
\foreach \x in {0.7,1.0,...,2.35}{\foreach \y in {0.4,0.7,1.0,1.3}{\fill[gray!60] (\x,\y) circle (0.04);}}
\draw[->,thick,blue!60!black] (G.west) -| (-0.5,0.9) -- (0.5,0.9);
\node[font=\scriptsize,blue!60!black,anchor=east] at (-0.6,2.1) {ball $\to$ current};
\draw[->,thick,red!70!black] (2.5,0.9) -- (3.5,0.9) |- (G.east);
\node[font=\scriptsize,red!70!black,anchor=west] at (3.6,2.1) {spikes $\to$ paddle};
\node[font=\scriptsize] at (1.5,-0.2) {neurons on an electrode array};
\node[font=\scriptsize,align=center] at (1.5,-0.85) {hit: predictable stimulus\\miss: random current};
% (b) reservoir
\begin{scope}[xshift=7.9cm]
\node[font=\small] at (1.6,4.3) {(b) reservoir};
\node[blk,text width=1.1cm] (X) at (-0.4,1.0) {input\\$u(t)$};
\draw[very thick,fill=orange!10] (1.6,1.0) circle (0.8);
\foreach \a/\r in {20/0.35,80/0.5,150/0.3,210/0.55,280/0.4,330/0.2,110/0.15}{\fill[red!60!black] ({1.6+\r*cos(\a)},{1.0+\r*sin(\a)}) circle (0.05);}
\node[font=\scriptsize] at (1.6,-0.2) {organoid: no teacher};
\node[blk,text width=1.8cm,fill=green!8] (W) at (4.0,1.0) {readout\\$W_{\text{out}}$};
\draw[->,thick] (X) -- (0.8,1.0);
\draw[->,thick] (2.4,1.0) -- node[above,font=\scriptsize] {$\mathbf r(t)$} (W);
\draw[->,thick] (W) -- (4.0,2.3) node[above,font=\scriptsize] {$y(t)$};
\node[font=\scriptsize,green!40!black] at (4.0,0.3) {supervised training};
\end{scope}
\end{tikzpicture}
\caption{Two ways to make a living network compute. (a) Closed loop: the ball position is delivered as current, spikes from the motor regions move the paddle, and a predictable or random stimulus after each hit or miss serves as the teacher. (b) Reservoir~\eqref{eq:reservoir}: the organoid gets no training signal; it splits the input into many nonlinear responses with memory, and a single linear readout outside learns from the error. The organoid's own connections change as well, without a teacher.}
\label{fig:livingmachine}
\end{figure}

\section{A living reservoir}

There is another way to use a living network: not to train it at all. In engineering this is called \emph{reservoir computing}~\cite{Maass2002,JaegerHaas2004}. Take a ready-made nonlinear system with memory, any system, as long as its response to the input is rich and depends on the recent past. The input $u(t)$ drives it, producing a response vector $\mathbf r(t)$, and only a linear readout is trained:
\begin{empheq}[box=\keybox]{equation}
y(t) \;=\; W_{\text{out}}\,\mathbf r(t) ,
\label{eq:reservoir}
\end{empheq}
with weights fitted by the least-mean-squares rule~\eqref{eq:lms} of Chapter~\ref{sec:motorlearning}. The reservoir does what the small \nt{GLUT} neurons of Chapter~\ref{sec:motorlearning} did: it splits the input into a long vector in which the desired classes become separable by a plane (Chapter~\ref{sec:dendrites}).

An organoid on an electrode array proved suitable as such a reservoir. Speech sounds delivered to it as current patterns were classified by speaker with about $78\%$ accuracy after training a single linear readout~\cite{Cai2023}. The $78\%$ figure is the one the authors report and the press repeats. The result is useful, but it matters where the ``intelligence'' in it lies: the organoid supplies nonlinearity and fading memory, while error-driven learning happens outside, in silicon (Fig.~\ref{fig:livingmachine}). The organoid does not stay unchanged, though: according to the authors, its own functional connectivity changed during training, and they call this unsupervised learning within the organoid itself~\cite{Cai2023}. How much of the accuracy comes from this change and how much from the readout has not been separated.

\section{What the bench tells us about the machine}

\emph{Energy.} By estimate~\eqref{eq:energy} a spike costs about $10^{-10}$~J. A culture of a million neurons firing at one spike per second spends on signaling
\begin{empheq}[box=\keybox]{equation}
P \;\approx\; 10^{6}\times1\,\text{s}^{-1}\times10^{-10}\,\text{J} \;=\; 10^{-4}\,\text{W},
\label{eq:dishpower}
\end{empheq}
a tenth of a milliwatt. The electronics that stimulate, record, and process these spikes spend orders of magnitude more. As in Chapter~\ref{sec:debugging}, the bottleneck is not the computation but the connection to it.

\emph{Missing parts.} The bench shows how much a \nt{GLUT} bus with \nt{GABA} flow control can do without control registers: self-organize, generate network bursts, change its responses under feedback, and serve as a good reservoir. What it cannot do without them is decide for itself what counts as success. On the bench that signal is supplied by the engineer; in the brain, as we saw in Chapters~\ref{sec:dofa}--\ref{sec:acet}, by the broadcast channels.

\emph{Ethics.} Organoids are grown from human cells, and the more complex they become, the more serious the question of whether such tissue can feel anything. The engineering view suggests a partial answer: pain in Chapter~\ref{sec:pain} is not the signal of a single sensor but a whole alarm system with gates, a volume control, and broadcast channels, none of which the bench has. But this is an argument, not a proof, and the field itself proposes carrying out ethical assessment alongside the experiments, from the very start~\cite{Smirnova2023}.

\begin{keyblock}
\begin{proposition}[A living network on the bench]
\label{prop:livingmachine}
A network of \nt{GLUT} and \nt{GABA} neurons on an electrode array generates network bursts on its own~\eqref{eq:burstinterval}, changes its responses under feedback, and can serve as a reservoir~\eqref{eq:reservoir} for a readout trained outside. All of this is measured. That such a network learns by some general rule, for example to predict its input, is a hypothesis, and most specialists do not accept loud claims about its ``sentience.''
\end{proposition}
\end{keyblock}

\section{Dopamine on a flash of light}
\label{sec:dishdopamine}

The only direct experiment we know of in which a culture was given dopamine as a teacher was run on a platform of organoids that researchers access remotely~\cite{Jordan2024}. Dopamine in a light-sensitive ``cage'' was added to the fluid bathing the organoids: the bound molecule does not act on receptors. The concentration of caged dopamine is $1$~mM. When the network is to be rewarded, it is illuminated with ultraviolet light: the cage breaks apart, and dopamine is released into the volume.

The loop works like this. The same stimulus is delivered over and over; if it evoked more spikes than before, a flash is triggered and dopamine is released. Over $240$ repetitions the number of evoked spikes grew overall. The authors themselves write that a single experiment of this kind cannot show that the growth was caused by dopamine~\cite{Jordan2024}.

For an engineer this is a first attempt to build the three-factor rule~\eqref{eq:threefactor} in a dish: sender and recipient activity leaves an eligibility trace, and a common signal decides whether to lock in the change. But here the signal can only be positive: reward is either present or absent, and the setup cannot deliver a negative error $\delta<0$. Dopamine also spreads out and is cleared slowly, like the concentration in~\eqref{eq:lowpass}, so frequent flashes may simply raise its overall level. To tell learning apart from this, two controls are needed: the same number of flashes at random times unrelated to the network's response, and the same flashes without caged dopamine. A test after rest is also needed: does the change persist once the dopamine is gone?

\section{Dopamine teaches silicon}
\label{sec:biohybrid}

In the reverse experiment, living cells teach electronics~\cite{Keene2020}. Dopamine-releasing cells are grown in a microchannel above an organic neuromorphic element, a transistor whose channel conductance plays the role of a synaptic weight. Dopamine released by the cells is oxidized at the element's electrode, and this changes the conductance. The device also mimics the mechanism that clears transmitter from the cleft, so the weight can not only be changed for a long time but also reset.

In circuit terms this is a leaky integrator with a chemical input: the weight accumulates the dopamine flux, and ``clearance'' sets how fast it is forgotten, the same RC circuit as in~\eqref{eq:lowpass}. What matters here is the direction of the link. The probe of Chapter~\ref{sec:debugging} cannot see the broadcast registers because they are chemical. This element reads exactly a chemical register and turns it into an electrical weight. For brain--computer interfaces this is a path toward a sensor that hears not only the data bus but also the control registers.

\section{Organoids with their own dopamine neurons}
\label{sec:midbrainorg}

Organoids resembling the brain region that houses the \nt{DOPA} neurons can be grown from human stem cells. They contain working dopamine neurons that release dopamine~\cite{Jo2016}. Such organoids are grown mainly to study disease. We found no experiments in which their dopamine served as a teacher for another network.

For a machine made of living neurons this is the missing part. The bench of Chapter~\ref{sec:livingmachine} is a data bus and flow control without control registers; here a source of the broadcast signal already exists. What is missing is a critic: a network that would compute the prediction error~\eqref{eq:ctd} and drive these neurons. Connecting a cortical organoid to such an organoid so that dopamine is released on error is an open problem, and for now it is a hypothesis, not an experiment.

\section{An organoid balances a pole without dopamine}
\label{sec:cartpole}

The most complete of the recent experiments does without dopamine altogether~\cite{Robbins2026}. Mouse cortical organoids were connected in a closed loop to the cart-pole task: organoid activity moves the cart, and the pole position is fed back by stimulation. Between trials the organoid receives short high-frequency training stimuli. Which ones is chosen by a reinforcement-learning program.

The results are measured:
\begin{itemize}
\item in most organoids, training signals chosen by the program gave better performance than randomly chosen signals or no signal;
\item after $45$ minutes of rest the improvement did not persist;
\item blocking both types of glutamate receptors, AMPA and NMDA, abolished the improvement.
\end{itemize}

The last point says where the learned change lives: in \nt{GLUT} synapses, and through the receptor that in Section~\ref{sec:weightstore} works as a detector of input--output coincidence. The second point shows what is missing: there is a fast change but no consolidation, like the fast learner of Chapter~\ref{sec:motorlearning} without the slow one. And the teacher's role has changed: the engineer choosing a current pattern has been replaced by a program, but the teacher still stands outside the dish.

Among earlier experiments on training cultures on electrode arrays we likewise found none that used dopamine: the training signal was always electrical stimulation.

\section{Who teaches the culture}

\begin{table}[t]
\centering
\footnotesize
\setlength{\tabcolsep}{4pt}
\renewcommand{\arraystretch}{1.15}
\begin{tabular}{>{\raggedright}p{2.7cm}>{\raggedright}p{2.5cm}>{\raggedright}p{3.2cm}>{\raggedright\arraybackslash}p{4.4cm}}
\toprule
Experiment & Who teaches & Training signal & What is known \\
\midrule
stimulation stops at the desired response & engineer & end of stimulation & the response is locked in: measured \\
DishBrain (ch.~\ref{sec:dishbrain}) & engineer & predictable or random current & significant growth of rallies within a session only with full feedback, a smaller effect with silence: measured; not robust from session to session; the cause is debated \\
cart-pole & reinforcement-learning program & high-frequency stimuli chosen by the program & better than random, but not retained after rest: measured \\
dopamine on a flash & rule ``more spikes means reward'' & dopamine released by light & growth in spikes: observation; role of dopamine not proven \\
biohybrid synapse & living cells & dopamine from cells & long-lasting change and reset of the element's weight: measured \\
organoids with \nt{DOPA} neurons & --- & --- & a dopamine source exists; not used as a teacher \\
\bottomrule
\end{tabular}
\caption{Who teaches living neurons on a chip, and with what.}
\label{tab:dishteachers}
\end{table}

In every experiment where the culture itself learns, the teacher is still outside (Table~\ref{tab:dishteachers}): an engineer, a program, or a rule set by an engineer. Dopamine has entered the dish only once, and its role remains to be proven. Building a real broadcast channel in a dish, with a critic, an error, and an eligibility trace as in Chapter~\ref{sec:dofa}, is the next step, and no one has taken it yet.


\section{Summary}

\begin{table}[t]
\centering
\footnotesize
\setlength{\tabcolsep}{4pt}
\renewcommand{\arraystretch}{1.15}
\begin{tabular}{>{\raggedright}p{2.8cm}>{\raggedright}p{4.2cm}>{\raggedright\arraybackslash}p{5.8cm}}
\toprule
Experiment & What the living network does & What is known \\
\midrule
network with no input & network bursts at intervals from a second to minutes~\eqref{eq:burstinterval} & measured; the mechanism via synaptic depletion is the accepted model \\
a teacher who falls silent & the desired response is locked in when stimulation is switched off & measured \\
playing tennis & runs of hits grow longer; significant improvement within a session only with full feedback & behavior change measured; learning from session to session not robust; effect size and the input-prediction explanation disputed \\
virtual body & movement toward a target with a chosen stimulus pattern & measured \\
reservoir & nonlinearity and memory for the readout~\eqref{eq:reservoir} & measured; error-driven learning is outside, in silicon; the organoid's connections also change \\
energy & about $10^{-4}$~W per million neurons~\eqref{eq:dishpower} & estimate from~\eqref{eq:energy} \\
\bottomrule
\end{tabular}
\caption{What machines made of living neurons have shown.}
\label{tab:livingmachine}
\end{table}

A machine made of living neurons (Table~\ref{tab:livingmachine}) is a bench that shows what a data bus and flow control can do without control registers. They can do a lot, but they cannot decide on their own what is good. On the bench that role is played by the engineer with a chosen current pattern; in the brain, by the few broadcast wires to which the middle of the book is devoted.

\section{Exercises}

\begin{exercise}[\lvlA]
\label{ex:21:1}
\emph{Interval between network bursts.} A network burst depletes the store to zero, then $\tau_{\text{rec}}\,\dd x/\dd t=1-x$, and a new network burst starts at $x=x^*$; $\tau_{\text{rec}}=0.8$~s. (a)~Find the interval for $x^*=0.9$ and $0.99$. (b)~The threshold $x^*$ varies randomly by $\pm0.005$ around $0.99$. What range do the intervals span?
\end{exercise}

\begin{exercise}[\lvlA]
\label{ex:21:2}
\emph{Energy of the bench.} (a)~What power does estimate~\eqref{eq:dishpower} give for the whole brain ($8.6\cdot10^{10}$ neurons)? (b)~The bench records $1024$ electrodes at $20$~kHz with $10$~bits per sample. At $1$~nJ per bit, how many times more does transmitting this stream cost than the power of the culture itself?
\end{exercise}

\begin{exercise}[\lvlB]
\label{ex:21:3}
\emph{Design: suppressing network bursts.} Sparse stimulation through many electrodes makes every synapse release transmitter at rate $r_b$; the store obeys $\tau_{\text{rec}}\,\dd x/\dd t=1-x-Ur_b\tau_{\text{rec}}x$, $U=0.5$, $\tau_{\text{rec}}=0.8$~s. What is the smallest $r_b$ at which the store never reaches $x^*=0.9$; $0.99$, and how much weaker does the synapse become?
\end{exercise}

\begin{exercise}[\lvlB]
\label{ex:21:4}
\emph{Reservoir memory is at most the number of neurons.} A reservoir with $N$ outputs $\mathbf r(t)$ receives a white input $u(t)$ with zero mean and unit variance; $m_k$ is the largest squared correlation of a readout $\mathbf w_k\cdot\mathbf r(t)$ with $u(t-k)$. Prove that the memory capacity $\mathrm{MC}=\sum_{k\ge0}m_k\le N$.
\end{exercise}

\begin{exercise}[\lvlB]
\label{ex:21:5}
\emph{Simulation: a reservoir in silicon.} Reservoir $\mathbf r(t+1)=\tanh(W\mathbf r(t)+\mathbf v\,u(t)+\mathbf c)$, $N=30$, $W$ random with spectral radius $0.9$, $v_i\sim U(-0.5,0.5)$, $c_i\sim U(-0.2,0.2)$, $u\sim U(-1,1)$, readout by ridge regression. Find the memory capacity of Exercise~\ref{ex:21:4} with and without $\tanh$. Which reservoir remembers more?
\end{exercise}

\begin{exercise}[\lvlB]
\label{ex:21:6}
\emph{A readout for a living reservoir.} An organoid provides $1024$ channels and $P=240$ training recordings of two classes. (a)~How many features $n$ can be kept so that, by the formula of Exercise~\ref{ex:19:3}, a random labeling is separable with probability at most $1\%$? (b)~How many sigmas above chance is an accuracy of $78\%$ on $100$ test recordings?
\end{exercise}

\begin{exercise}[\lvlC]
\label{ex:21:7}
\emph{Research: a teacher without a broadcast channel.} Propose a stimulation protocol that implements the three-factor rule of Chapter~\ref{sec:dofa} on the bench through $8$--$1024$ electrodes, and a frozen-readout control that distinguishes weight learning from a shift in network state. What result would refute the learning hypothesis?
\end{exercise}
